\input{../tex-inputs/latex-leseansicht-vorspann}

\section[Arthur Schnitzler an Gustav Schwarzkopf, 19. 8. {[}1913?{]}]{L04443 Arthur Schnitzler an Gustav Schwarzkopf, 19. 8. {[}1913?{]}}
\nopagebreak\mylabel{L04443v}
\rehead{ }\normalsize\beginnumbering\briefempfaengerindex{Schwarzkopf, Gustav@\textsc{Schwarzkopf, Gustav}!zzzSchnitzler, Arthur@\emph{von Arthur Schnitzler}!1913-08-191@{19. 8. {[}1913?{]}}|(be}
\toendnotes[C]{\smallbreak\pagebreak[2]}
\correspDesc{Versand  durch Arthur Schnitzler am 19. 8. [1913?] in Brijuni
\newline{}Übermittlung  am 19. 8. 1913 in Brijuni
\newline{}Erhalt  durch Gustav Schwarzkopf im Zeitraum [20. 8. 1913
                  – 24. 8. 1913?] in Wien}\toendnotes[C]{\smallbreak}
\Standort{DLA, A:Schnitzler, HS.1985.1.1897.}
\physDesc{Bildpostkarte, 391 Zeichen
\newline{}Handschrift: Bleistift, deutsche Kurrent
\newline{}Versand: Stempel: »\nobreak{}\oindex{Brijuni@\textbf{Brijuni}|pwk}Bri\textcolor{gray}{o}{[}ni{]}, 19. 8. 1\textcolor{gray}{3}\nobreak{}«.  }\toendnotes[C]{\smallbreak}\pstart{}{\pb}Hrn \textsc{Gustav Schwarzkopf}\pend{}\pstart{}\textsc{Wien I}\oindex{I., Innere Stadt@\textbf{I., Innere Stadt}|pw}\pend{}\pstart{}\textsc{Tiefer Graben 17}\oindex{Wien@\textbf{Wien}!I., Innere Stadt@\textbf{I., Innere Stadt}!Tiefer Graben 17@\textbf{Tiefer Graben 17}|pw}\pend{}{\bigskip}
\pstart
           \noindent{}\centering{}{\pb}\textcolor{gray}{\textbf{Insel BRIONI in der Adria\oindex{Brijuni@\textbf{Brijuni}|pw}}}\pend
           
\pstart
           \centering{}\textcolor{gray}{\textbf{Blick in eine \begin{otherlanguage}{italian}Cava Romana\end{otherlanguage}.}}\pend
           \vspace{1em}
\pstart
           \raggedleft{}{\pb}\textsc{Brioni}\oindex{Brijuni@\textbf{Brijuni}|pw}{ }19/8.\pend
           \vspace{0.5em}
\pstart
           lieber Guſtav,{ }\label{K_L04443-1v}\edtext{am Freitag}{\lemma{\textnormal{\emph{am Freitag}}}\Cendnote{\textnormal{Der \emph{Tagebuch}\pwindex{Schnitzler, Arthur 15. 5. 1862 Wien – 21. 10. 1931 ebd.@\textsc{Schnitzler, Arthur} (15. 5. 1862 Wien – 21. 10. 1931 ebd.)!Tagebuch@\strich\emph{Tagebuch}|pwk}eintrag von Freitag, dem 22. 8. 1913, bestätigt die Angaben der Karte (die
                  Kinder\pwindex{Schnitzler, Heinrich 9.\,8.\,1902 Hinterbrühl – 12.\,7.\,1982 Wien@\textsc{Schnitzler, Heinrich} (9.\,8.\,1902 Hinterbrühl – 12.\,7.\,1982 Wien)|pwkv}\pwindex{Cappellini, Lili 13.\,9.\,1909 Wien – 26.\,7.\,1928 Venedig@\textsc{Cappellini, Lili} (13.\,9.\,1909 Wien – 26.\,7.\,1928 Venedig)|pwkv} reisten nach Wien\oindex{Wien@\textbf{Wien}|pwk}, die
                  Eltern\pwindex{Schnitzler, Olga 17.\,1.\,1882 Wien – 13.\,1.\,1970 Lugano@\textsc{Schnitzler, Olga} (17.\,1.\,1882 Wien – 13.\,1.\,1970 Lugano)|pwkv} mit dem Schiff über Nacht nach Venedig\oindex{Venedig@\textbf{Venedig}|pwk}) und damit die Datierung des
                  Poststempels.}}}\label{K_L04443-1} fahren wir nach Venedig\oindex{Venedig@\textbf{Venedig}|pw},
                  \textsc{Grand Hotel}\oindex{Grand Hotel Venezia@\textbf{Grand Hotel Venezia}|pw} am So{\geminationn}tag, je nach Wetter Engadin\oindex{Engadin@\textbf{Engadin}|pw} oder oberital Seen\oindex{Oberitalienische Seen@\textbf{Oberitalienische Seen}|pw} (ev. Florenz\oindex{Florenz@\textbf{Florenz}|pw}); die Kinder\pwindex{Schnitzler, Heinrich 9.\,8.\,1902 Hinterbrühl – 12.\,7.\,1982 Wien@\textsc{Schnitzler, Heinrich} (9.\,8.\,1902 Hinterbrühl – 12.\,7.\,1982 Wien)|pwv}\pwindex{Cappellini, Lili 13.\,9.\,1909 Wien – 26.\,7.\,1928 Venedig@\textsc{Cappellini, Lili} (13.\,9.\,1909 Wien – 26.\,7.\,1928 Venedig)|pwv} direkt nach
                  Wien\oindex{Wien@\textbf{Wien}|pw}. We{\geminationn} es Ihnen
               Vergnügen macht, beſuchen Sie sie, ja?\pend
           
\pstart
           Eine Zeile über Ihr Befinden {\pb}nach Venedig\oindex{Venedig@\textbf{Venedig}|pw} würde uns \uline{ſehr} freuen.
               Heuer war es fortdauernd{ }ſehr ſchön.\pend
           
\pstart
           Herzlichſt mit Grüßen von uns Allen Ihr {\\[\baselineskip]}\spacefill\mbox{Arth}\pend
           \leftskip=0em{}\selectlanguage{ngerman}\endnumbering\briefempfaengerindex{Schwarzkopf, Gustav@\textsc{Schwarzkopf, Gustav}!zzzSchnitzler, Arthur@\emph{von Arthur Schnitzler}!1913-08-191@{19. 8. {[}1913?{]}}|)be}\mylabel{L04443h}
\begin{anhang}
\end{anhang}\newcommand{\dateiname}{L04443}\newcommand{\titel}{Arthur Schnitzler an Gustav Schwarzkopf, 19. 8. [1913?]}\newcommand{\editorInnen}{Herausgegeben von Jahnke, SelmaMüller, Martin Anton}\input{../tex-inputs/latex-leseansicht-abspann}
